
\subsection{\texorpdfstring{\textbf{
      Динамика классической и непрерывной модели Хопфилда}}
  {Динамика классической и непрерывной модели Хопфилда}}

\textbf{Для дискретного случая} и бинарного паттерна с компонентами из
\(\{ 0,1\}\) имеем уравнение вида:

\[
  x_{i} = \ \left\{ \begin{array}{r}
    1\ if\ \sum_{j \neq \ i}^{}{T_{ij}x_{j}} \geq \theta_{i},\ \ \\
    0\ otherwise.
  \end{array} \right.\
\]

или так (\(T_{ii}\  = \ 0\))


\[
  x_{i} = \ \left\{ \begin{array}{r}
    1\ if\ \sum_{}^{}{{\langle T}_{\{ row = i\}},x\rangle} \geq \theta_{i},\ \ \\
    0\ otherwise.
  \end{array} \right.\
\]


где \(\theta_{i}\) - погрешность активации, параметр модели. \cite{hopfield1982neural}

Причём, можно заменить \(\{ 0,1\}\) на \(\{ - 1,1\}\) в зависимости от
правил обучения.

Динамика может быть асинхронная (произвольная активация или наложен
порядок активации) и синхронная (одновременная активация - нереалистично).

\textbf{Для непрерывного случая} имеем дифференциальное уравнение вида:

\[
  \tau_{f}\ \frac{dx_{i}}{dt} = \ \sum_{j = 1}^{N_{f}}{T_{ij}V_{j}}\  - \ x_{i}\  + \ I_{i}
\]


T - матрица весов связей, \(V_{i}\) = функция активации
\(g(x_{i}) = \frac{\partial L_{x}}{\partial x_{i}}\) \cite{hopfield1984neurons}

Есть также современная динамика с использованием скрытых слоёв, но здесь
мы её не рассматриваем. \cite{ramsauer2020hopfield}
